\appendix


\addcontentsline{toc}{chapter}{\protect\numberline{} \vspace{0pt}\hspace*{-0.3in} Appendix}

In Chapter \ref{chap:2nd}, we exhibit the combinatorial proof of Corollary \ref{cor:nocore}, and this corollary can also be proved by using analytic method. In a similar manner, we have the following proposition.

\begin{sprop}
Let $p_d(n)$ be the number of partitions of n into distinct parts. For $|q|<1$,
$$\sum_{n=0}^{\infty}  p_d(n) q^n =\left(\sum_{n=0}^{\infty} p(n) q^{2n}\right) \left(\sum_{n=0}^{\infty} q^{\frac{n(n+1)}{2}}\right).$$
\end{sprop}

\vspace{3mm}
We guess that it is possible to find a combinatorial proof of this proposition.

\vspace{3mm}
In Chapters \ref{chap:3rd} and \ref{chap:4th}, we enumerate $(s,s+d,\dots,s+pd)$-core partitions and self-conjugate $(s,s+1,\dots,s+p)$-core partitions. It would be very interesting if one can find a lattice path interpretation or closed formula for self-conjugate $(s,s+d,\dots,s+pd)$-core partitions.